\chapter{Formulación de Galerkin y Ecuación de Poisson}
\label{ap:galerkin-poisson}

Este apéndice desarrolla con rigor las dos piezas formales sobre las que se apoya el Capítulo \ref{ch:bsplines-galerkin}. Primero, presentamos el método de Galerkin como caso particular de los métodos de residuos ponderados y lo aplicamos a la discretización de la ecuación radial de Hartree-Fock. Después, mostramos cómo el mismo formalismo permite tratar la interacción electrónica: definimos el potencial radial $Y^k(a,b;r)$ a partir de la expansión multipolar del operador de Coulomb y demostramos, mediante diferenciación bajo el signo integral (regla de Leibniz generalizada), que dicho potencial satisface una ecuación de Poisson radial generalizada.

\section{Formulación de Galerkin}
\label{sec:galerkin-formulacion}

\subsection{La formulación fuerte}

Empezamos con la ecuación de Schrödinger radial para un solo electrón en un potencial central:

\[ - \frac{1}{2}\dv[2]{P}{r} + V_{\text{eff}}(r) P(r) = \epsilon P(r), \]

donde \( V_{\text{eff}}(r) \) incluye la atracción nuclear, el término centrífugo \( \frac{l(l+1)}{2r^2} \), y cualquier potencial \( SCF \) de interés (Hartree/Intercambio).

\subsection{Expansión en la base}

Aproximamos la función de onda radial \( P(r) \) como una combinación lineal de B-splines \( B_j(r) \):

\begin{equation}
\label{eq:basis-exp}
  P(r) \approx \sum_{j=1}^N c_jB_j(r),
\end{equation}

\begin{itemize}
\item \( c_j \) son los coeficientes que necesitamos encontrar.
\item \( N \) es la cantidad de splines.
\end{itemize}

\subsection{Método de residuos ponderados}

Como la expansión en B-splines es una aproximación, al sustituirla en la ecuación de Schrödinger no se obtiene cero exactamente; se obtiene un residuo \( R(r) \). El objetivo es minimizar este residuo proyectando la ecuación sobre una función de peso (o prueba) \( w_i(r) \) y forzando a que la integral ponderada sea nula sobre el dominio \( [0,R_{\text{max}}] \). Existen diversas estrategias para seleccionar estas funciones de peso, lo que da lugar a diferentes formulaciones numéricas, como se resume en \cref{tab:mwr-methods}.

\begin{table}[htbp]
    \centering
    \small
    \caption{Comparación de Métodos de Residuos Ponderados (MWR) para la minimización del residuo $R$.}
    \label{tab:mwr-methods}
    \renewcommand{\arraystretch}{1.6}
    \setlength{\tabcolsep}{4pt}

    \begin{tabularx}{\textwidth}{
        >{\RaggedRight\bfseries}m{2.1cm}
        >{\RaggedRight}m{3.0cm}
        >{\centering}m{2.9cm}
        >{\RaggedRight\arraybackslash}X
    }
        \toprule
        Método & Función Peso ($w_i$) & Condición & Interpretación \newline Geométrica \\
        \midrule

        Galerkin &
        Funciones base \newline $w_i = B_i$ &
        $\int R \cdot B_i \, dr = 0$ &
        El \textbf{residuo} es ortogonal al espacio de las funciones base. Proyección óptima en términos de energía. \\
        \cmidrule{1-4}

        Colocación \newline \textit{\footnotesize (Collocation)} &
        Delta de Dirac \newline $w_i = \delta(r - r_i)$ &
        $R(r_i) = 0$ &
        Fuerza a que el \textbf{residuo se anule} estrictamente en los nodos de control seleccionados. \\
        \cmidrule{1-4}

        Mínimos Cuadrados \newline \textit{\footnotesize (Least \newline Squares)} &
        Derivada del \newline residuo \newline $w_i = \frac{\partial R}{\partial c_i}$ &
        $\int R \cdot \frac{\partial R}{\partial c_i} \, dr = 0$ &
        Minimiza la norma $L^2$ del \textbf{residuo global} ($\int R^2 dr$). Ideal para operadores no simétricos. \\
        \cmidrule{1-4}

        Subdominio \newline \textit{\footnotesize (Finite \newline Volume)} &
        Función escalón \newline $w_i = \begin{cases} 1 & r \in \Omega_i \\ 0 & \text{resto} \end{cases}$ &
        $\int_{\Omega_i} R \, dr = 0$ &
        El \textbf{promedio integral} del residuo es cero dentro de cada volumen de control $\Omega_i$. \\
        \bottomrule
    \end{tabularx}
\end{table}

En este trabajo utilizamos el método de Galerkin, donde las funciones de peso son idénticas a las funciones base (\( w_i = B_i \)). Como se menciona en la tabla, esto implica que el residuo es ortogonal al espacio generado por las funciones base.

Esencialmente estamos buscando la mejor aproximación posible a la solución real dentro del espacio limitado generado por las funciones base, i.e. B-splines. Así, al forzar que el error residual sea ortogonal a cada una de estas funciones, aseguramos geométricamente que no quede ninguna componente del error ``proyectada'' dentro del subespacio que podemos controlar. Esto es análogo a minimizar la distancia entre un punto y un plano trazando una línea perpendicular: se elimina cualquier parte del error que pudiera haber sido corregida ajustando los coeficientes. De esta forma, garantizamos que la discrepancia restante sea puramente intrínseca a las limitaciones del tamaño de la base elegida y no a la calidad de nuestra proyección.

\subsection{Forma débil}

Multiplicamos la ecuación entera por una función base arbitraria \( B_i(r) \) e integramos sobre el dominio \( [0,R_{\text{max}}] \)

\begin{equation}
\label{eq:mwr1}
  \int_{0}^{R_{\text{max}}} B_i(r) \qty[ -\frac{1}{2}\dv[2]{r} + V_{\text{eff}}(r) - \epsilon ] P(r) \dd{r} = 0.
\end{equation}

Al reacomodar términos se obtiene

\[ \underbrace{\int_{0}^{R_{\text{max}}} B_i(r) \left( -\frac{1}{2} \frac{d^2}{dr^2} \right) P(r) \, dr}_{\text{Término Cinético}} + \underbrace{\int_{0}^{R_{\text{max}}} B_i(r) V_{\text{eff}}(r) P(r) \, dr}_{\text{Término Potencial}} = \epsilon \underbrace{\int_{0}^{R_{\text{max}}} B_i(r) P(r) \, dr}_{\text{Término Traslape}}. \]

El término cinético presenta una segunda derivada que se puede tratar mediante integración por partes. Esto no solo reduce el órden de derivación requerido para las funciones base, i.e. permite el uso de bases de menor continuidad, sino que también simetriza el operador, lo cual simplifica el álgebra lineal posterior.

\begin{equation} \label{eq:k-ibp}
  \int_{0}^{R_{\text{max}}} B_i(r) P''(r) \dd{r} = \underbrace{\eval{ B_i(r) P'(r) }_{0}^{R_{\text{max}}}}_{\text{Término de borde}} - \int_{0}^{R_{\text{max}}} B_i'(r) P'(r) \dd{r}.
\end{equation}

El término de borde se anula por las condiciones de frontera sobre la función de onda \( P(0) = P(R_{\text{max}})=0 \).

\subsection{Ecuación matricial}

Sustituimos la expansión (\cref{eq:basis-exp}) en las ecuaciones integrales. El intercambio entre la sumatoria y el signo integral no requiere aquí ninguna justificación de convergencia: la expansión es una suma \textit{finita} de $N$ términos, y la linealidad de la integral basta para factorizarla. La cuestión de la convergencia uniforme, que sí sería necesaria para una serie infinita, se traslada en este esquema a la pregunta distinta de con qué rapidez la solución discreta se aproxima a la exacta al aumentar $N$.

\begin{itemize}
\item La integral cinética
  \[ \frac{1}{2} \int \dv{B_i}{r} \qty(\sum_jc_j\dv{B_{j}}{r})\dd{r} = \sum_jc_j\underbrace{\qty[\frac{1}{2}\int \dv{B_i}{r}\dv{B_j}{r}\dd{r}]}_{T_{ij}}. \]
\item La integral potencial
  \[ \int B_i(r)V_{eff}(r) \qty(\sum_jc_jB_j(r)) \dd{r} = \sum_jc_j \underbrace{\qty[\int B_i(r)V_{\text{eff}}(r)B_j(r)\dd{r}]}_{V_{ij}}. \]
\item La integral de traslape
  \[ \epsilon \int B_i(r) \qty(\sum_jc_jB_j(r)) \dd{r} = \epsilon \sum_jc_j \underbrace{\qty[\int B_i(r)B_j(r)\dd{r}]}_{S_{ij}}. \]
\end{itemize}

Al reunir todos las términos para todo \( i = 1 \dots N \), obtenemos un sistema de ecuaciones lineales:

\[ \sum_j (T_{ij} + V_{ij})c_j = \epsilon \sum_j S_{ij}c_j. \]

Definimos el elemento de matriz Hamiltoniana como \( H_{ij} = T_{ij} + V_{ij} \) y el sistema de ecuaciones se convierte al problema de eigenvalores generalizado:

\begin{equation}
\label{eq:eigen}
\bm{H}\bm{c} = \epsilon \bm{S}\bm{c}
\end{equation}

\subsection{Equivalencia formal entre los métodos de Ritz y Galerkin}
\label{sec:ritz-galerkin}

En los textos estándar de química cuántica se utiliza el principio variacional de Ritz, pero en este trabajo utilizamos el método de Galerkin para la discretización de las ecuaciones. En el contexto de operadores Hermitianos, como el Hamiltoniano atómico, esta elección no altera la física subyacente, dado que ambos enfoques convergen a la misma estructura algebraica.

El método de Ritz parte de la búsqueda de puntos estacionarios del funcional de energía para una función de onda de prueba $|\tilde{\psi}\rangle$:

\begin{equation}
    E[\tilde{\psi}] = \frac{\langle \tilde{\psi} | \hat{H} | \tilde{\psi} \rangle}{\langle \tilde{\psi} | \tilde{\psi} \rangle}
\end{equation}

Al introducir una base finita $|\tilde{\psi}\rangle = \displaystyle\sum_{j} c_j |B_j\rangle$, la minimización con respecto a los coeficientes ($\pdv{E}{c_i^{*}} = 0$) conduce a la condición de que la proyección del operador desplazado sobre cada función base sea nula:

\begin{equation}
    \label{eq:ritz_condition}
    \sum_{j} c_j \langle B_i | \hat{H} - E | B_j \rangle = 0
\end{equation}

Por otra parte, el método de Galerkin aborda el problema desde la teoría de residuos ponderados. Se define el vector residual $|R\rangle$ como la discrepancia generada al aplicar el operador a la aproximación:

\begin{equation}
    |R\rangle = (\hat{H} - E) |\tilde{\psi}\rangle
\end{equation}

La condición estricta de Galerkin impone que el residuo sea ortogonal al subespacio generado por las funciones base. Esto es, proyectar el residuo sobre cada bra \( \langle B_i| \):

\begin{equation}
    \langle B_i | R \rangle = 0
\end{equation}

Sustituyendo la expansión lineal de $|\tilde{\psi}\rangle$ en esta proyección, obtenemos:

\begin{equation}
  \mel{B_i}{(\hat{H} - E)}{\sum_jc_jB_j} = 0
\end{equation}

Aprovechando la linealidad del producto interno, factorizamos los coeficientes y la suma, recuperando una expresión idéntica a (\cref{eq:ritz_condition}) derivada mediante Ritz:

\begin{equation}
    \sum_{j} c_j \left( \langle B_i | \hat{H} | B_j \rangle - E \langle B_i | B_j \rangle \right) = 0
\end{equation}

Esta identidad confirma que la proyección de Galerkin es matemáticamente equivalente a la variación de Ritz para este sistema. En consecuencia, ambos métodos conducen al mismo problema generalizado de valores propios:

\begin{equation}
    \mathbf{H}\mathbf{c} = E \mathbf{S}\mathbf{c}
\end{equation}

donde $\mathbf{H}$ y $\mathbf{S}$ son las matrices Hamiltoniana y de traslape respectivamente.

La equivalencia justifica la elección que hacemos en este trabajo. Si ambos caminos conducen al mismo problema matricial, conviene tomar el que resulte más directo de implementar, y ese es el de Galerkin. La vía de Ritz exige construir el funcional de energía, derivarlo respecto de cada coeficiente e igualar a cero, lo que en la práctica significa manipular simbólicamente una expresión cuadrática antes de poder escribir una sola línea de código. La proyección de Galerkin, en cambio, entrega los elementos de matriz directamente como integrales de productos de funciones base, que es exactamente la forma que la cuadratura numérica consume. La diferencia no es de rigor ni de resultado, sino de distancia entre la formulación y el algoritmo.

\section{Interacción Electrónica y el Potencial \texorpdfstring{$Y^k$}{Yk}}
\label{sec:interaccion-electronica}

En un sistema atómico con \( N \) electrones y un núcleo fijo de carga \( Z \), partimos del Hamiltoniano no relativista en unidades atómicas ya derivado en el Apéndice \ref{ap:unidades-atomicas} (Ec. \eqref{eq:hamiltonian1}):

\[
  H = \sum_i^N \qty(- \frac{1}{2}\nabla_i^2 - \frac{Z}{r_i}) + \sum_{i<j}^N \frac{1}{\abs{\bm{r}_i-\bm{r}_j}}.
\]

El término que presenta mayor dificultad computacional y teórica es la repulsión interelectrónica, representada por el operador de Coulomb \( V_{ij} = \abs{\bm{r}_i-\bm{r}_j}^{-1} \). Debido a que este operador acopla las coordenadas de los pares electrónicos, impide separar directamente las variables en la ecuación de Schrödinger. Para tratar este término en una geometría esférica, se utiliza la expansión de Laplace del potencial de Coulomb como armónicos esféricos.

\subsection{Expansión multipolar del potencial de Coulomb}

Consideremos dos electrones \( i \) y \( j \) con vectores de posición \( r_i \) y \( r_j \). De acuerdo con el teorema de adición de armónicos esféricos, cuya formulación en el lenguaje de operadores tensoriales irreducibles se desarrolla en el Apéndice \ref{ap:operadores-tensoriales}, la interacción recíproca puede expandirse como:

\begin{equation}
  \label{eq:laplace1}
  \frac{1}{\abs{\bm{r}_i - \bm{r}_j}} = \sum_{k=0}^{\infty} \frac{4\pi}{2k+1} \frac{r_<^k}{r_>^{k+1}} \sum_{q = -k}^k Y_{kq}^{*}(\Omega_i) Y_{kq}(\Omega_j),
\end{equation}

donde:

\begin{itemize}
\item \( r_< = \text{mín}(r_i,r_j) \) y \( r_> = \text{máx}(r_i,r_j) \).
\item \( \Omega_i = (\theta_i, \phi_i) \) denota las coordenadas angulares del electrón \( i \).
\item \( Y_{kq} \) son los armónicos esféricos normalizados.
\end{itemize}

Esta expansión permite desacoplar la dependenncia radial de la angular. La parte angular se resuelve analíticamente mediante el álgebra de momento angular, mientras que la física de la interacción de apantallamiento queda contenida enteramente en la parte radial.

\subsection{Potenciales radiales \texorpdfstring{$Y^k$}{Yk}}

En el formalismo de Hartree-Fock, la evaluación de la energía y de las ecuaciones de Fock requiere calcular la interacción electrostática entre pares de orbitales. Independientemente de si tratamos con la interacción directa o de intercambio, el término fundamental que describe la distribución espacial de los electrones es el producto de dos funciones de onda radiales:

\begin{equation}
\label{eq:pair-density}
    \rho_{ab}(r) = P_a(r) P_b(r).
\end{equation}

Es crucial notar que esta cantidad $\rho_{ab}(r)$ posee una interpretación física profunda: representa la proyección radial de los elementos de la \textbf{matriz de densidad de una partícula}, $\gamma(\mathbf{r}, \mathbf{r}')$, en el espacio de configuraciones.

\begin{itemize}
\item Cuando $a = b$, $\rho_{aa}(r)$ corresponde a la densidad de carga electrónica clásica de la capa $a$.
\item Cuando $a \neq b$, $\rho_{ab}(r)$ representa una \textit{densidad de transición} o de intercambio. Aunque no tiene un análogo en la electrostática clásica, matemáticamente actúa como una distribución de carga ``virtual'' que genera un campo de potencial efectivo.
\end{itemize}

El componente radial del potencial electrostático de orden multipolar $k$ generado por esta distribución $\rho_{ab}(s)$, actuando sobre un electrón en la posición $r$, se define formalmente mediante la integración sobre la coordenada $s$ de la distribución fuente:

\begin{equation}
 \label{eq:vk-integral}
 v_k(a,b;r) = \int_0^{\infty} \frac{r_<^k}{r_>^{k+1}} P_a(s)P_b(s)\dd{s}.
\end{equation}

Para evaluar esta integral, es necesario dividir el intervalo de integración $[0,\infty)$ en dos regiones, dependiendo de la posición relativa de la variable de integración $s$ respecto al punto de observación $r$:

\begin{itemize}
\item \textbf{Región interior} ($s < r$): En este caso, $r_< = s$ y $r_> = r$. El integrando contribuye con un factor de $s^k/r^{k+1}$, representando el efecto de apantallamiento de la densidad interna.
\item \textbf{Región exterior} ($s > r$): En este caso, $r_< = r$ y $r_> = s$. El integrando contribuye con un factor $r^k/s^{k+1}$, representando el potencial multipolar generado por la densidad externa.
\end{itemize}

Sustituyendo estas condiciones en \cref{eq:vk-integral}, obtenemos la expresión explícita:

\begin{equation}
 \label{eq:vk-integral2}
 v_k(a, b; r) = \frac{1}{r^{k+1}} \int_0^r s^k P_a(s) P_b(s) \, ds + r^k \int_r^{\infty} \frac{1}{s^{k+1}} P_a(s) P_b(s) \, ds.
\end{equation}

Siguiendo el formalismo de Froese Fischer \cite{fischer1977}, resulta conveniente trabajar con una función de potencial escalada por $r$, lo cual simplifica la estructura de la ecuación radial de Schrödinger (eliminando singularidades y factores de $1/r$ en el término de potencial). Definimos así las funciones $Y^k$:

\begin{equation}
 \label{eq:yk-pot}
 Y^k(a,b;r) \equiv r \cdot v_k(a,b;r).
\end{equation}

Al aplicar este factor $r$ a la ecuación (\cref{eq:vk-integral2}), obtenemos la forma integral definitiva que constituye el núcleo de nuestro tratamiento numérico:

\begin{equation}
\label{eq:Yk-final}
 Y^k(a,b;r) = r^{-k}\int_0^r s^kP_a(s)P_b(s) \dd{s} + r^{k+1} \int_r^{\infty} s^{-(k+1)}P_a(s)P_b(s)\dd{s}.
\end{equation}

Esta función $Y^k(a,b;r)$ encapsula toda la información necesaria para construir tanto los operadores de Coulomb $J$ (donde $a=b$) como los operadores de Intercambio $K$ (donde $a \neq b$) para cualquier átomo multielectrónico bajo la aproximación de campo central. Permite además expresar las integrales de Slater de interacción directa ($F^k$) e intercambio ($G^k$) como simples valores esperados radiales, por ejemplo:

\begin{equation}
\label{eq:Fk-def}
    F^k(nl, n'l') = \int_0^{\infty} P_{nl}(r) \frac{Y^k(n'l', n'l'; r)}{r} P_{nl}(r) \dd{r}.
\end{equation}

Si bien \cref{eq:Yk-final} proporciona la definición formal y exacta del potencial, su evaluación numérica directa mediante cuadraturas presenta desventajas computacionales significativas. La naturaleza global del operador integral implica que el valor del potencial en un punto depende de la función en todo el dominio, lo que en una discretización conduce a matrices densas y operaciones de costo $O(N^2)$.

Sin embargo, como se demuestra en adeltante, mediante la aplicación de la regla de Leibniz, la función $Y^k(a,b;r)$ es solución de una ecuación diferencial de segundo orden con condiciones de frontera específicas. Esta ecuación, conocida como la \textit{Ecuación de Poisson Radial Generalizada}, toma la forma:

\begin{equation}
\label{eq:poisson-radial}
    \left( -\frac{d^2}{dr^2} + \frac{k(k+1)}{r^2} \right) Y^k(a,b;r) = \frac{2k+1}{r} P_a(r)P_b(r).
\end{equation}

Para completar la definición del problema de valores de contorno se imponen las condiciones que se deducen del comportamiento de la integral original, y que se derivan en detalle en la Sección \ref{sec:condiciones-frontera}:
\begin{equation}
    Y^k(a,b; 0) = 0, \qquad Y^k(a,b; r \to \infty) \longrightarrow \frac{M_k}{r^{k}},
\end{equation}
donde $M_k$ es el momento multipolar de orden $k$ de la densidad. Nótese que el decaimiento de $Y^k$ va como $r^{-k}$ y no como $r^{-(k+1)}$, que es el del potencial sin escalar $v_k = Y^k/r$.

La transición del formalismo integral al diferencial (\cref{eq:poisson-radial}) ofrece ventajas decisivas para la implementación numérica, particularmente en el contexto del Método de Elementos Finitos (FEM) que empleamos en este trabajo:

\begin{enumerate}
    \item \textbf{Localidad y Esparsidad:} A diferencia del operador integral, el operador diferencial es local. Al proyectar esta ecuación sobre una base de B-splines (o cualquier base de soporte compacto), las matrices resultantes son de banda estrecha (dispersas). Esto reduce la complejidad computacional de la inversión matricial y el almacenamiento de memoria de $O(N^2)$ a $O(N)$.

    \item \textbf{Consistencia con el Solucionador:} Dado que las funciones de onda radiales $P(r)$ se obtienen resolviendo la ecuación de Schrödinger (una ecuación diferencial de segundo orden), tratar los potenciales $Y^k$ bajo el mismo esquema diferencial permite unificar la arquitectura del \textit{solver}. Tanto los orbitales como los potenciales se resuelven mediante el mismo mecanismo de Galerkin descrito en la Sección \ref{sec:galerkin-formulacion}.

    \item \textbf{Tratamiento de Singularidades:} La Ecuación (\cref{eq:poisson-radial}) permite tratar de manera robusta el término centrífugo $k(k+1)/r^2$ y la singularidad de Coulomb $1/r$ en el origen mediante las propiedades de las funciones base, evitando las inestabilidades numéricas que suelen surgir al integrar numéricamente kernels singulares.
\end{enumerate}

En consecuencia, en este trabajo adoptamos (\cref{eq:poisson-radial}) como la ecuación gobernante para determinar los potenciales de interacción, transformando el problema de Hartree-Fock en un sistema acoplado de ecuaciones diferenciales. Su discretización por el método de Galerkin, que sigue paso por paso el procedimiento de la Sección \ref{sec:galerkin-formulacion} aplicado ahora al operador de Poisson, se desarrolla en la Sección \ref{subsec:poisson-metodo} del capítulo de método numérico, donde conduce a la forma débil \eqref{eq:poisson-weak-metodo} y al sistema lineal \eqref{eq:poisson-linear-metodo}; no la repetimos aquí. A continuación demostramos formalmente la equivalencia entre la forma diferencial y la definición integral (\cref{eq:Yk-final}), que es lo que este apéndice aporta y que el capítulo da por establecido.

\section{Diferenciación bajo el signo integral}
\label{sec:leibniz}

Para fundamentar la demostración del teorema principal de esta sección, enunciamos primero dos resultados clásicos del cálculo integral que servirán como herramientas auxiliares. Estos teoremas se presentan sin demostración, asumiendo su validez en el contexto de las funciones continuas tratadas en este trabajo.

\begin{theorem}[Teorema del valor medio para integrales]\label{thm:valor-medio}
  Si \( f \) es una función continua en un intervalo cerrado \( [a,b] \), entonces existe al menos un punto \( c \in [a,b] \) tal que:
  \[
  \int_{a}^{b} f(x) \dd{x} = f(c)(b-a).
  \]
\end{theorem}

\begin{theorem}[Regla de Leibniz con límites constantes]\label{thm:leibniz-constante}
  Sea \( f(s,r) \) una función continua sobre un rectángulo \( R = [a,b] \times [c,d] \) tal que la derivada parcial \( \pdv{f}{r} \) existe y es continua. Entonces:
  \[
  \dv{r} \int_{a}^{b} f(s,r) \dd{s} = \int_{a}^{b} \pdv{f}{r}(s,r) \dd{s}.
  \]
\end{theorem}

A continuación, generalizamos estos conceptos para el caso donde los límites de integración son funciones de la variable de diferenciación.

\begin{theorem}[Regla general de Leibniz]\label{thm:leibniz-general}

  Considérese una función \( f(s,r) \) definida y continua sobre un rectángulo \( R = [a,b]\times [c,d] \subset \mathbb{R}^2 \). Supongamos que se cumplen las siguientes condiciones:

  \begin{enumerate}
  \item\label{item:1} La derivada parcial \( \pdv{f}{r} \) existe y es continua en el dominio de interés.
  \item\label{item:2} Las funciones de los límites \( \alpha, \beta : [c,d] \to [a,b] \) son diferenciables en su intervalo de definición.
  \end{enumerate}

  En tal caso, la función paramétrica \( F:[c,d] \to \mathbb{R} \) dada por

  \[ F(r) = \int_{\alpha(r)}^{\beta(r)} f(s,r) \dd{s} \]

  es diferenciable para todo \( r \in (c,d) \), y su derivada está dada por:

  \begin{equation}
    \label{eq:leibniz-int}
   \dv{F}{r} = \int_{\alpha(r)}^{\beta(r)} \pdv{f}{r}(s,r) \dd{s} + f(\beta(r),r)\beta'(r) - f(\alpha(r),r)\alpha'(r).
  \end{equation}

\end{theorem}

\begin{proof}

  Sea \( r_0 \in (c,d) \). Para determinar la derivada de \( F \) en \( r_0 \), consideramos el cociente incremental:

  \[ \frac{F(r) - F(r_0)}{r-r_0} = \frac{1}{r-r_0} \qty(\int_{\alpha(r)}^{\beta(r)} f(s,r) \dd{s} - \int_{\alpha(r_0)}^{\beta(r_0)} f(s,r_0) \dd{s} ). \]

  Mediante la aditividad del dominio respecto a los puntos \( \alpha(r_0) \) y \( \beta(r_0) \), podemos descomponer la primera integral de la siguiente manera:

  \[ \int_{\alpha(r)}^{\beta(r)} f(s,r) \dd{s} = \int_{\alpha(r)}^{\alpha(r_0)}f(s,r) \dd{s} + \int_{\alpha(r_0)}^{\beta(r_0)} f(s,r)\dd{s} + \int_{\beta(r_0)}^{\beta(r)} f(s,r)\dd{s}. \]

  Al sustituir esta expresión en el cociente anterior y agrupar términos, obtenemos tres componentes que analizaremos por separado:

  \[
  \begin{aligned}
  \frac{F(r) - F(r_0)}{r - r_0} &= \underbrace{\frac{1}{r - r_0} \int_{\alpha(r)}^{\alpha(r_0)} f(s,r) \dd{s}}_{I_1} \\
  &\quad + \underbrace{\int_{\alpha(r_0)}^{\beta(r_0)} \frac{f(s,r) - f(s,r_0)}{r - r_0} \dd{s}}_{I_2} \\
  &\quad + \underbrace{\frac{1}{r - r_0} \int_{\beta(r_0)}^{\beta(r)} f(s,r) \dd{s}}_{I_3}.
  \end{aligned}
  \]

  A continuación, calculamos el límite de cada término cuando \( r \to r_0 \):

  \textbf{1. Términos de borde (\( I_1 \) y \( I_3 \)):} De acuerdo con el teorema del valor medio para integrales (\cref{thm:valor-medio}), existe un valor \( \bar{s}_1 \) entre \( \alpha(r) \) y \( \alpha(r_0) \) tal que:

  \[ I_1 = \frac{\alpha(r_0) - \alpha(r)}{r-r_0} f(\bar{s}_1,r) = - \frac{\alpha(r) - \alpha(r_0)}{r-r_0}f(\bar{s}_1,r). \]

  Cuando \( r \to r_0 \), sabemos que \( \bar{s}_1 \to \alpha(r_0) \). Por la definición de derivada y la continuidad de \( f \), obtenemos:

  \[ \lim_{r\to r_0} I_{1} = - f(\alpha(r_0), r_0) \alpha'(r_0). \]

  De forma análoga para \( I_3 \), existe \( \bar{s}_2 \) entre \( \beta(r_0) \) y \( \beta(r) \), lo que resulta en:

  \[ \lim_{r\to r_0} I_3 = f(\beta(r_0),r_0)\beta'(r_0). \]

  \textbf{2. Término integral (\( I_2 \)):} Dado que \( \pdv{f}{r} \) es continua en el rectángulo compacto \( R \), por la regla de Leibniz para límites constantes (\cref{thm:leibniz-constante}) es posible intercambiar el límite y la integral:

  \[ \lim_{r \to r_0} I_2 = \int_{\alpha(r_0)}^{\beta(r_0)} \lim_{r\to r_0} \frac{f(s,r) - f(s,r_0)}{r-r_0} \dd{s} = \int_{\alpha(r_0)}^{\beta(r_0)} \pdv{f}{r}(s,r_0) \dd{s}. \]

  Finalmente, sumando los tres límites resultantes, concluimos que \( F \) es derivable en \( r_0 \) y su derivada es:

  \[ F'(r_0) = \int_{\alpha(r_0)}^{\beta(r_0)} \pdv{r}f(s,r_0) \dd{s} + f(\beta(r_0),r_0) \beta'(r_0) - f(\alpha(r_0), r_0) \alpha'(r_0). \]

  Como \( r_0 \) representa un punto arbitrario del intervalo \( (c,d) \), la igualdad se sostiene para todo el dominio. Esto establece la validez de la \cref{eq:leibniz-int}.

\end{proof}

\section{Derivadas del Potencial \texorpdfstring{$Y^k$}{Yk}}
\label{sec:derivadas-yk}

Partimos de la definición integral escalada del potencial radial (donde definimos a \( v_k \) como el potencial integral asintótico o bruto, y \( Y^k = r v_k \) como el potencial de Hartree-Fock radial escalado) (\cref{eq:Yk-final}):

\begin{equation}
    Y^k(r) = r^{-k} \mathcal{I}_{\text{int}}(r) + r^{k+1} \mathcal{I}_{\text{ext}}(r),
\end{equation}

donde definimos las integrales auxiliares como:

\begin{align}
    \mathcal{I}_{\text{int}}(r) &= \int_{0}^{r} s^k \rho_{ab}(s) \, ds, \\
    \mathcal{I}_{\text{ext}}(r) &= \int_{r}^{\infty} s^{-(k+1)} \rho_{ab}(s) \, ds.
\end{align}

Aplicamos la regla de Leibniz (\cref{thm:leibniz-general}) para encontrar las derivadas de estas integrales auxiliares:

\begin{itemize}
    \item Para $\mathcal{I}_{\text{int}}(r)$: $a(r)=0$, $b(r)=r$.
    \begin{equation}
        \frac{d}{dr}\mathcal{I}_{\text{int}} = (r^k \rho_{ab}(r)) \cdot (1) - (0) \cdot (0) = r^k \rho_{ab}(r).
    \end{equation}

    \item Para $\mathcal{I}_{\text{ext}}(r)$: $a(r)=r$, $b(r)=\infty$.
    \begin{equation}
        \frac{d}{dr}\mathcal{I}_{\text{ext}} = (0) \cdot (0) - (r^{-(k+1)} \rho_{ab}(r)) \cdot (1) = -r^{-(k+1)} \rho_{ab}(r).
    \end{equation}
\end{itemize}

En ambos casos hemos usado que el integrando depende solo de la variable de integración y no del parámetro $r$. La regla de Leibniz para un integrando $f(s)$ con límites variables consta de tres términos, y el que involucra la derivada del integrando se anula de manera idéntica:

\[ \dv{r}\int_{a(r)}^{b(r)} f(s)\dd{s} = f(b(r))\,b'(r) - f(a(r))\,a'(r) + \int_{a(r)}^{b(r)} \pdv{f(s)}{r} \dd{s}, \qquad \pdv{f(s)}{r} = 0. \]

Solo sobreviven, por tanto, las contribuciones de los límites de integración, que son las que se han evaluado arriba.

Ahora, derivamos $Y^k(r)$ utilizando la regla del producto y sustituyendo los resultados anteriores:

\begin{align*}
    \frac{dY^k}{dr} &= \frac{d}{dr}\left( r^{-k} \mathcal{I}_{\text{int}} \right) + \frac{d}{dr}\left( r^{k+1} \mathcal{I}_{\text{ext}} \right) \\
    &= \left[ -k r^{-k-1} \mathcal{I}_{\text{int}} + r^{-k} \left( \frac{d\mathcal{I}_{\text{int}}}{dr} \right) \right] + \left[ (k+1) r^k \mathcal{I}_{\text{ext}} + r^{k+1} \left( \frac{d\mathcal{I}_{\text{ext}}}{dr} \right) \right] \\
    &= -k r^{-k-1} \mathcal{I}_{\text{int}} + r^{-k} (r^k \rho_{ab}) + (k+1) r^k \mathcal{I}_{\text{ext}} + r^{k+1} (-r^{-(k+1)} \rho_{ab}) \\
    &= -k r^{-k-1} \mathcal{I}_{\text{int}} + \rho_{ab}(r) + (k+1) r^k \mathcal{I}_{\text{ext}} - \rho_{ab}(r).
\end{align*}

Observamos que los términos de fuente $\rho_{ab}(r)$ se cancelan exactamente, lo que confirma la continuidad del campo. La primera derivada resultante es:

\begin{equation}
\label{eq:Yk-prime-app}
    \frac{dY^k}{dr} = -k r^{-k-1} \mathcal{I}_{\text{int}}(r) + (k+1) r^k \mathcal{I}_{\text{ext}}(r).
\end{equation}

\subsubsection{Segunda Derivada y Ecuación de Poisson}

Diferenciamos \cref{eq:Yk-prime-app} nuevamente respecto a $r$:

\begin{align*}
    \frac{d^2Y^k}{dr^2} &= \frac{d}{dr}\left( -k r^{-k-1} \mathcal{I}_{\text{int}} \right) + \frac{d}{dr}\left( (k+1) r^k \mathcal{I}_{\text{ext}} \right) \\
    &= \left[ k(k+1) r^{-k-2} \mathcal{I}_{\text{int}} - k r^{-k-1} \left( \frac{d\mathcal{I}_{\text{int}}}{dr} \right) \right] \\
    &\quad + \left[ k(k+1) r^{k-1} \mathcal{I}_{\text{ext}} + (k+1) r^k \left( \frac{d\mathcal{I}_{\text{ext}}}{dr} \right) \right].
\end{align*}

Sustituimos nuevamente las derivadas de las integrales auxiliares ($\mathcal{I}'_{\text{int}} = r^k \rho_{ab}$ y $\mathcal{I}'_{\text{ext}} = -r^{-(k+1)} \rho_{ab}$):

\begin{align*}
    \frac{d^2Y^k}{dr^2} &= k(k+1) r^{-k-2} \mathcal{I}_{\text{int}} - k r^{-k-1} (r^k \rho_{ab}) \\
    &\quad + k(k+1) r^{k-1} \mathcal{I}_{\text{ext}} + (k+1) r^k (-r^{-(k+1)} \rho_{ab}).
\end{align*}

Factorizamos el término común $\frac{k(k+1)}{r^2}$ en los términos que contienen integrales y agrupamos los términos que contienen $\rho_{ab}$:

\begin{align*}
    \frac{d^2Y^k}{dr^2} &= \frac{k(k+1)}{r^2} \left( \underbrace{r^{-k} \mathcal{I}_{\text{int}} + r^{k+1} \mathcal{I}_{\text{ext}}}_{Y^k(r)} \right) - \frac{k}{r} \rho_{ab} - \frac{k+1}{r} \rho_{ab} \\
    &= \frac{k(k+1)}{r^2} Y^k(r) - \frac{2k+1}{r} \rho_{ab}(r).
\end{align*}

Finalmente, recuperamos la ecuación diferencial de segundo orden utilizada en el método de Galerkin:

\begin{equation}
  \label{eq:poisson-true}
    \frac{d^2 Y^k}{dr^2} = \frac{k(k+1)}{r^2} Y^k(r) - \frac{2k+1}{r} P_a(r)P_b(r).
\end{equation}

\subsubsection{Análisis de las Condiciones de Frontera}
\label{sec:condiciones-frontera}

Para que la ecuación diferencial de segundo orden derivada anteriormente tenga una solución única, es necesario especificar las condiciones de frontera en los límites del dominio radial, es decir, en $r \to 0$ y $r \to \infty$. Estas condiciones emergen naturalmente del comportamiento asintótico de la definición integral (\cref{eq:Yk-final}).

\paragraph{Comportamiento en el origen ($r \to 0$).}

Consideremos el comportamiento de los dos términos de $Y^k(r)$ cuando $r$ es pequeño. Sabemos que las funciones de onda radiales $P(r)$ se comportan cerca del origen como $r^{l+1}$, por lo que la densidad de par se comporta como $\rho_{ab}(r) \sim r^{l_a + l_b + 2}$.

Analizamos el primer término (contribución interna):
\begin{equation}
    \lim_{r \to 0} \left[ r^{-k} \int_0^r s^k \rho_{ab}(s) \, ds \right].
\end{equation}
Dado que el intervalo de integración tiende a cero, la integral se anula más rápido que la singularidad $r^{-k}$ (asumiendo $k$ dentro de los límites físicos permitidos por los momentos angulares $l_a, l_b$). Específicamente, el integrando es de orden $s^{k + l_a + l_b + 2}$, por lo que la integral escala como $r^{k + l_a + l_b + 3}$. El término completo se comporta como $r^{l_a + l_b + 3}$, que tiende a 0.

Analizamos el segundo término (contribución externa):
\begin{equation}
    \lim_{r \to 0} \left[ r^{k+1} \int_r^{\infty} s^{-(k+1)} \rho_{ab}(s) \, ds \right].
\end{equation}
La integral converge a un valor finito constante (el momento multipolar total del átomo). El factor $r^{k+1}$ domina, llevando el término a 0 (dado que $k \geq 0$).

Por lo tanto, la condición de frontera Dirichlet en el origen es estricta:
\begin{equation}
\label{eq:bc-origin}
    Y^k(a,b; 0) = 0.
\end{equation}

\paragraph{Comportamiento asintótico ($r \to \infty$)}

Cuando $r$ es muy grande (fuera de la extensión espacial de la nube electrónica), la densidad $\rho_{ab}(r)$ decae exponencialmente a cero. Esto tiene dos consecuencias sobre las integrales auxiliares:

1.  $\mathcal{I}_{\text{ext}}(r) = \int_r^\infty s^{-(k+1)}\rho_{ab}(s) ds \to 0$ rápidamente.
2.  $\mathcal{I}_{\text{int}}(r) = \int_0^r s^k \rho_{ab}(s) ds$ converge a una constante $M_k$, que representa el momento multipolar eléctrico de orden $k$ de la distribución de carga.

Sustituyendo esto en la ecuación de $Y^k$:
\begin{equation}
    Y^k(r \to \infty) \approx r^{-k} M_k + r^{k+1} (0) = \frac{M_k}{r^k}.
\end{equation}

Esta relación nos proporciona la condición de frontera adecuada para un dominio computacional truncado en un radio finito $R_{\text{max}}$ suficientemente grande. Dependiendo del valor de $k$:

\begin{itemize}
    \item \textbf{Caso $k=0$ (Monopolo):} Si $a=b$ (densidad de carga del orbital), $M_0 = \int_0^\infty P_a^2 ds = 1$ (por normalización). Entonces $Y^0(r) \to 1$. Esto representa el apantallamiento total de la carga nuclear vista desde el infinito.
    \item \textbf{Caso $k > 0$:} El potencial decae a cero.
\end{itemize}

Escrita como condición de Dirichlet, la prescripción exterior sería
\begin{equation}
\label{eq:bc-final}
    Y^k(0) = 0 \quad \text{y} \quad Y^k(R_{\text{max}}) = \frac{1}{R_{\text{max}}^k} \int_0^{R_{\text{max}}} s^k P_a(s) P_b(s) \, ds,
\end{equation}
y en la práctica es común aproximarla para $k \neq 0$ por $Y^k(R_{\text{max}}) = 0$ cuando $R_{\text{max}}$ es suficientemente grande. Ambas formas tienen el inconveniente de exigir el momento multipolar como dato de entrada, siendo que depende de la densidad que se está calculando.

La alternativa que adoptamos en la implementación evita ese acoplamiento. Derivando la relación asintótica $Y^k \to M_k r^{-k}$ se elimina la constante desconocida y queda una condición homogénea sobre el borde,
\begin{equation}
\label{eq:bc-robin}
    \left. \dv{Y^k}{r} \right|_{R_{\text{max}}} + \frac{k}{R_{\text{max}}}\, Y^k(R_{\text{max}}) = 0,
\end{equation}
que es de tipo Robin y que en la formulación débil se implementa sumando $k/R_{\text{max}}$ al último elemento diagonal de la matriz del operador, sin alterar su simetría ni su carácter definido positivo. Para el monopolo se reduce a la condición natural de Neumann, que reproduce el límite constante $Y^0 \to M_0$ dejando que la solución determine el valor de la meseta. La condición en el origen, por su parte, no se impone como ecuación sino eliminando del espacio de prueba los primeros $k+1$ splines, que son los únicos incompatibles con el comportamiento $Y^k \sim r^{k+1}$.

%%% Local Variables:
%%% mode: LaTeX
%%% TeX-master: "../main"
%%% End:
